\scanpage{48}
\mspage{46}
% archive sheet number: 47
On trouve :

\uline{Proposition} 9.4. \inserted{a)} Pour \deleted{\unsure{établir}} la validité du th. de monodromie \inserted{pour les $H^i_!(X_{\overline K},\unsure{\mathbb Q_\ell|Q_\ell})$}, pour $V$ fixé, \deleted{et} des $X_K$ avec $\dim X_K\leq n$, et tout $i$, il suffit de savoir qu'il est vrai pour tout $X_{K'}$, où $X$ est lisse \inserted{de dim $\leq n$} sur \deleted{$\unsure{K'}$} \inserted{une} extension $K'$ finie radicale de $K$, considéré comme corps des fractions de la clôture normale $V'$ de $V$ dans $K'$. \deleted{\unsure{Et}} On peut \deleted{supposer que les $X_{K'}$} se borner à supposer de plus les $X_{K'}$ affines, \deleted{\unsure{et}} \inserted{b)} Dans le cas où \deleted{la résolution des singularités si} toute extension de t.f. de degré de transcendance $\leq n$ sur la clôture parfaite $K^{\unsure{p^{-\infty}|p^\infty}}$ de $K$ admet un modèle projectif et lisse, il suffit de supposer \inserted{les} $X_{K'}$ projectifs. c) Enfin, \deleted{si} les conditions \deleted{\unsure{précédentes}} \inserted{\deleted{\unsure{ci-dessus}}} de a) sont vérifiées, alors pour $\forall$ $X_K$ de t.f. et \uline{lisse} sur $K$, \inserted{de dim $\leq n$,} le th. de monodromie est vrai aussi pour les $H^i_{\deleted{!}}(X_{\overline K},\mathbb Q_\ell)$.

9.5. Nous allons maintenant examiner le th de monodromie pour $H^i(X_{\overline K},\mathbb Q_\ell)$, \deleted{\unsure{donc}} \inserted{pour} $X_K$ de t.f., en supposant que \deleted{les conditions de 9.4. a) sont satisfaites \unsure{pour tout $n$|et tout $n$}. On procède par récurrence sur $i$, le cas $i<0$ étant trivial,} c'est prouvé pour les \deleted{\unsure{les}} $H^{\unsure{j|i}}(X_{\overline K},\mathbb Q_\ell)$ pour $j<i$.
